%% kapa: exam=1 round=33 year=2022 period=2 form=A subject=accounting subject_ko=회계학 seq=27 number=67
%% source: Q-Net 감정평가사 기출문제 (cst00309) artlSeq=5209997 file=raw/1008/5209997/2022년도 제33회 감정평가사 1차시험 2교시(감정평가 관계법규, 회계학).pdf
%% conversion: kordoc 4.24.51 -> tools/gen1.py (auto); review=proofread-sonnet; flags=none
\begin{question}{67}
(주)감평은 20×1년 9월 1일 미국에 있는 토지(유형자산)를 \$5,000에 취득하고 원가모형을 적용하고 있다. 20×1년 12월 31일 현재 토지의 공정가치는 \$5,100이며, 20×2년 2월 1일 토지 중 30\%를 \$1,550에 처분하였다. 일자별 환율이 다음과 같을 때, 처분손익은? (단, (주)감평의 기능통화는 원화이다.)
\par{\small\centering\begin{tabular}{|c|c|c|c|}\hline
일 자 & 20×1년 9월 1일 & 20×1년 12월 31일 & 20×2년 2월 1일 \\ \hline
환율(￦/\$) & ￦1,200 & ￦1,170 & ￦1,180 \\ \hline
\end{tabular}\par}
\begin{choices}
  \item 손실 ￦29,000
  \item 손실 ￦38,900
  \item ￦0
  \item 이익 ￦29,000
  \item 이익 ￦38,900
\end{choices}
\end{question}
\answer{4}
